\chapter{Introdução}

A promulgação da Lei de Acesso à Informação (LAI nº 12.527/2011) representou um marco na democratização da gestão pública brasileira, consolidando o direito constitucional do cidadão em acessar informações produzidas pelo Estado \cite{Lei:12.527:2011}. Como desdobramento dessa legislação, os portais de transparência municipais tornaram-se instrumentos obrigatórios para publicação de dados orçamentários, financeiros e administrativos. No entanto, a mera disponibilização de dados brutos não garante, por si só, o exercício efetivo do controle social \cite{matiaspereira2010}. A complexidade técnica da linguagem orçamentária, a fragmentação das informações em múltiplas tabelas e relatórios, e a ausência de ferramentas analíticas acessíveis criam barreiras significativas entre o cidadão e a compreensão real de como os recursos públicos são aplicados \cite{melo2019proposta}. Nesse contexto, o \textit{Data Warehouse} emerge como uma solução tecnológica capaz de integrar, organizar e transformar dados dispersos em informação estratégica e inteligível, viabilizando análises que transcendem a simples consulta de valores isolados \cite{kimball2013data}.

A modelagem dimensional, especialmente através da metodologia consolidada por Ralph Kimball e Margy Ross, oferece uma arquitetura orientada à usabilidade analítica, diferentemente dos modelos relacionais tradicionais focados em otimização transacional \cite{kimball2013data}. Ao estruturar dados orçamentários em esquemas dimensionais --- com tabelas de fatos representando transações financeiras e tabelas de dimensões contextualizando essas operações por tempo, órgão, função governamental, favorecido e outros atributos ---, o \textit{Data Warehouse} permite que perguntas complexas sejam respondidas de forma intuitiva: ``Quanto foi investido em educação nos últimos três anos?'', ``Qual o tempo médio de pagamento a fornecedores na área de saúde?'', ``Existe concentração excessiva de recursos em poucos favorecidos?''. Essas questões, fundamentais para o controle social, tornam-se acessíveis não apenas a especialistas em finanças públicas, mas a cidadãos, conselheiros municipais, jornalistas e organizações da sociedade civil \cite{marco2022transparencia}.

A transparência orçamentária no Brasil, embora amparada por marcos legais robustos, enfrenta desafios significativos na transição entre disponibilidade formal de dados e efetividade do controle social. Estudos demonstram que a simples publicação de informações em portais governamentais não resulta automaticamente em \textit{accountability} ou participação cidadã qualificada. A literatura aponta três lacunas críticas que justificam a relevância desta pesquisa: a lacuna social, referente às barreiras de literacia financeira e digital que impedem a população de compreender e questionar a aplicação de recursos públicos; a lacuna técnica, relacionada à fragmentação, inconsistência e baixa usabilidade dos dados disponibilizados; e a lacuna pedagógica, que diz respeito à ausência de instrumentos que traduzam informação técnica em conhecimento acessível e acionável.

Do ponto de vista social, pesquisas sobre controle social evidenciam que a participação cidadã é constrangida não apenas pela falta de acesso à informação, mas fundamentalmente pela dificuldade em interpretá-la e conectá-la a consequências tangíveis na vida cotidiana. O orçamento público, apresentado em categorias contábeis como ``elemento de despesa'', ``função programática'' e ``fonte de recurso'', permanece distante da realidade experimentada pelos cidadãos nas escolas, postos de saúde e infraestrutura urbana. Essa opacidade técnica desestimula a participação e perpetua assimetrias de poder: apenas aqueles com formação especializada ou acesso a intermediários (como ONGs e jornalistas investigativos) conseguem exercer fiscalização efetiva. A consequência é um déficit democrático onde a transparência formal coexiste com a opacidade prática.

Sob a perspectiva técnica, os portais de transparência municipais frequentemente disponibilizam dados em formatos inadequados para análise (PDFs não estruturados, planilhas sem padronização, ausência de APIs), dificultam o cruzamento de informações entre diferentes bases (despesas, receitas, licitações) e carecem de recursos analíticos que permitam visualizações temporais, comparações multidimensionais ou identificação de padrões anômalos. A ausência de uma arquitetura de dados integrada obriga usuários avançados a realizar processos manuais e repetitivos de coleta, limpeza e consolidação --- um investimento de tempo e especialização que a maioria dos cidadãos não pode fazer. O \textit{Data Warehouse}, ao consolidar esses processos em uma estrutura permanente e reutilizável, democratiza não apenas o acesso, mas a capacidade analítica em si.

A dimensão pedagógica desta justificativa encontra fundamento nos princípios freireanos de educação libertadora e conscientização. Paulo Freire argumentava que a transformação social depende da capacidade dos oprimidos de ``ler o mundo'' criticamente, superando a ``cultura do silêncio'' imposta por estruturas de dominação \cite{freire1967liberdade}. Aplicado ao contexto orçamentário, isso significa que a democratização do poder requer mais do que transparência passiva: exige ferramentas que capacitem o cidadão a formular perguntas, testar hipóteses e construir narrativas próprias sobre a gestão pública. O conceito freireano de ``inédito viável'' --- aquilo que é possível mas ainda não realizado --- aplica-se perfeitamente ao \textit{Data Warehouse} como tecnologia que viabiliza uma forma ainda incipiente de participação: a análise pública orientada por dados, onde cidadãos podem questionar não apenas ``quanto foi gasto'', mas ``por que essa área recebeu mais recursos'', ``esse investimento está alinhado com as necessidades locais'' ou ``há indícios de favorecimento a determinados fornecedores'' \cite{freire1996autonomia}.

A escolha da execução orçamentária municipal como objeto de estudo reforça a relevância social do trabalho: é na esfera local que as políticas públicas materializam-se mais concretamente no cotidiano dos cidadãos, e é também onde as possibilidades de controle social direto são maiores \cite{novaes2014orcamento}. Ao desenvolver um modelo replicável de \textit{Data Warehouse} para gestão municipal, este TCC contribui não apenas para a transparência de uma cidade específica, mas oferece um referencial metodológico que pode ser adaptado por outros municípios, ampliando o alcance da iniciativa. Ademais, a integração entre fundamentos teóricos de controle social (ciências sociais aplicadas) e modelagem dimensional (ciência da computação) preenche uma lacuna interdisciplinar: a maior parte dos trabalhos sobre transparência pública permanece no campo da análise política ou jurídica, enquanto estudos sobre \textit{Data Warehouse} raramente abordam suas implicações cívicas. Esta pesquisa posiciona-se na intersecção dessas áreas, argumentando que a arquitetura de dados não é neutra --- ela é uma escolha política que define quem pode participar do debate público e em que condições.

\section{Objetivos}

\subsection*{Objetivo Geral}

Desenvolver um \textit{Data Warehouse} para análise da execução orçamentária do município de Juiz de Fora, reorganizando dados públicos dispersos em estrutura dimensional acessível que facilite a interpretação e fiscalização cidadã dos gastos públicos.

\subsection*{Objetivos Específicos}

\begin{itemize}
    \item Coletar e documentar dados de receitas e despesas do Portal da Transparência da Prefeitura de Juiz de Fora, identificando inconsistências nas fontes disponíveis.
    \item Modelar dimensionalmente a execução orçamentária segundo metodologia Kimball, mapeando classificações oficiais (MCASP/MTO) em fatos e dimensões.
    \item Implementar processos de ETL para extração, limpeza, transformação e carga dos dados, tratando heterogeneidades de formato e nomenclatura.
    \item Construir \textit{dashboards} interativos com indicadores como execução orçamentária por área, gasto per capita e concentração de fornecedores.
    \item Avaliar a usabilidade das visualizações produzidas em comparação ao formato original dos dados, refletindo sobre o papel da arquitetura de dados na democratização da informação pública.
\end{itemize}
